\begin{abcliste}
\abc
Von $0$ bis $t_1$ wirkt eine konstante Kraft, also ist die Beschleunigung konstant (zweites Newtonsches Gesetz): Die Geschwindigkeit steigt gleichmässig, eine steigende Gerade. Nach $t_1$ wirkt keine horizontale Kraft mehr (Gewichtskraft und Normalkraft heben sich auf, keine Reibung): Nach dem ersten Newtonschen Gesetz bleibt die Geschwindigkeit konstant, eine waagrechte Linie.
\begin{center}
\begin{tikzpicture}[>=Stealth, xscale=0.8, yscale=1.6]
 \draw[->] (0,0) -- (8.4,0) node[right] {$t$ in \si{\second}};
 \draw[->] (0,0) -- (0,1.75) node[above] {$v$ in \si{\mps}};
 \foreach \x in {1,...,8} \draw (\x,0.03) -- (\x,-0.03) node[below] {\small \x};
 \foreach \y in {0.5,1.0,1.5} \draw (0.08,\y) -- (-0.08,\y) node[left] {\small \y};
 \draw[ultra thick] (0,0) -- (3.6,1.38) -- (7.5,1.38);
 \draw[dashed] (3.6,0) -- (3.6,1.38);
 \node[below] at (3.6,-0.18) {$t_1$};
 \draw[dashed] (7.5,0) -- (7.5,1.38);
 \node[below] at (7.5,-0.18) {$t_2$};
\end{tikzpicture}
\end{center}

\abc
Ab $t_1$ bleibt die Geschwindigkeit gleich:
\[
 v(t_2) = v_1 = \result{\vEO}
\]
Der Weg ist die Fläche unter dem Graphen, ein Dreieck und ein Rechteck:
\begin{align}
 s &= \sZF \\
   &= \frac12\cdot\vE\cdot\tE+\vE\cdot(\tZ-\tE) \\
   &= \sZ \approx \result{\sZP{0}{2}}
\end{align}

\abc
Die Beschleunigung von $0$ bis $t_1$ ist $a = v_1/t_1 = \aWP{0}{2}$, und die Kraft
\begin{align}
 F &= m\,a = \FPF \\
   &= \mW\cdot\frac{\vE}{\tE} \\
   &= \FP \approx \result{\FPP{0}{2}}
\end{align}

\abc
Lea irrt sich. Der Propeller hat dem Wagen keinen „Schwung“ als Kraftvorrat mitgegeben, der sich danach aufbraucht. Eine Kraft ändert die Geschwindigkeit; sie ist nicht nötig, um eine Geschwindigkeit zu erhalten. Nach $t_1$ wirkt keine horizontale Kraft, also \result{\text{fährt der Wagen mit konstanter Geschwindigkeit weiter}}. Langsamer würde er nur, wenn eine Kraft gegen die Bewegung wirkt, etwa Reibung oder Luftwiderstand, und die sind hier vernachlässigbar.
\end{abcliste}
